\section{Configuração da Medição de Energia}

\subsection{Instrumentação e Procedimento}
\textbf{ESP32 e Raspberry Pi:} corrente e tensão foram amostradas a cada 500 ms usando o INA219. Duas janelas foram capturadas: repouso ($N_{idle}$ amostras) e carga ($N_{load}$ amostras) durante a execução do benchmark.\newline
\textbf{Dispositivos Android:} corrente e tensão foram lidas das APIs do sistema sob fases controladas de repouso/carga.
A instrumentação completa e as janelas de amostragem estão resumidas na Tabela~\ref{tab:energy_setup}.

A potência incremental foi estimada como $\Delta P = P_{load} - P_{idle}$, e a energia por inferência como $E_{inf} = \Delta P / \text{vazão}$.

\rev{A Tabela~\ref{tab:android_specs} detalha os três dispositivos Android (A--C) usados na campanha. Como a inferência executa na CPU por meio do interpretador TFLite (XNNPACK, sem delegate de GPU/NNAPI), a configuração dos clusters de CPU e a memória disponível são os recursos relevantes para a carga de IA; as GPUs são listadas apenas por completude.}

\begin{table}[t]
\caption{\rev{Especificações dos dispositivos Android (dados do fabricante; a RAM entre parênteses é o MemTotal reportado pelo sistema).}}
\label{tab:android_specs}
\centering
\scriptsize
\setlength{\tabcolsep}{3pt}
\begin{tabular}{lccc}
\toprule
 & Android A & Android B & Android C \\
\midrule
Aparelho & Redmi Note 12 & Galaxy A56 5G & Redmi Note 11S \\
Modelo & 23021RAAEG & SM-A566E & 2201117SG \\
SoC & Snapdragon 685 & Exynos 1580 & Helio G96 \\
Litografia & 6 nm & 4 nm & 12 nm \\
CPU & 4$\times$A73 2,8 & 1$\times$A720 2,9 & 2$\times$A76 2,05 \\
(GHz) & 4$\times$A53 1,9 & 3$\times$A720 2,6 & 6$\times$A55 2,0 \\
 & & 4$\times$A520 1,95 & \\
GPU & Adreno 610 & Xclipse 540 & Mali-G57 MC2 \\
RAM & 6 GB (5,77) & 8 GB (7,60) & 6 GB (5,86) \\
Android & 15 & 16 & 13 \\
\bottomrule
\end{tabular}
\end{table}

\begin{table}[t]
\caption{Configuração medida e janelas de amostragem por plataforma.}
\label{tab:energy_setup}
\centering
\footnotesize
\setlength{\tabcolsep}{4pt}
\begin{tabular}{lccc}
\toprule
Plataforma & Fonte & $N_{idle}$ & $N_{load}$ \\
\midrule
ESP32 (MicroPy) & INA219 & 50 & 50 \\
Raspberry Pi 4B & INA219 & 97 & 440 \\
Android (A-C) & Telem. sistema & 200 & 200 \\
\bottomrule
\end{tabular}
\end{table}


O envelope de implantação resultante combina vazão e energia incremental por inferência: cenários de alta taxa favorecem a execução em smartphones e SBCs, enquanto o sensoriamento autônomo de baixa taxa ainda pode ser suportado em hardware da classe microcontrolador.
